\subsection{理性化命题的证明细节}

\begin{proposition}[理性化与最大元]
若 $\succeq$ 完备且传递，且 $B$ 有限非空，则
$C(B)=\{x\in B:x\succeq y,\ \forall y\in B\}$ 非空。
\end{proposition}
\begin{proof}
从 $B$ 中任取 $x_1$。若它击败所有元素，则完成；否则存在 $x_2$ 使 $x_2\succ x_1$。继续比较。由于 $B$ 有限，严格改进不可能无限进行，过程在某个 $x_k$ 停止。若仍存在 $y\succ x_k$，应继续改进，矛盾。因此 $x_k$ 是最大元，且属于选择集。
\end{proof}

\begin{example}[并列最大元]
若 $a\sim b\succ c$，则 $C(\{a,b,c\})=\{a,b\}$。理性化并不意味着选择对应必须单值；只有严格偏好或严格拟凹等附加条件才能保证唯一选择。
\end{example}

\begin{proposition}[严格偏好环排除理性化]
若有限选择数据的显示偏好闭包包含严格环，则不存在完备传递偏好关系理性化该数据。
\end{proposition}
\begin{proof}
设严格环为 $x^1\succ x^2\succ\cdots\succ x^k\succ x^1$。严格偏好的传递性推出 $x^1\succ x^1$，违反严格偏好的反自反性。因此任何理性偏好都不可能包含该环。
\end{proof}

\subsection{从选择数据到效用不等式}
若 $x^t$ 在预算 $B^t$ 中被选择，且 $x^s$ 也可负担，则理性化要求 $u(x^t)\ge u(x^s)$。一串可负担关系便产生一串效用弱不等式；若最后回到起点并至少有一次严格不可负担，则得到 $u(x^t)>u(x^t)$ 的矛盾。这是 GARP 的算术核心。

\begin{example}[一条三期不等式链]
若 $x^1$ 在第 1 期可负担 $x^2$，$x^2$ 在第 2 期可负担 $x^3$，而第 3 期严格排除 $x^1$，则理性化要求
$u(x^1)\ge u(x^2)\ge u(x^3)>u(x^1)$，显然不可能。该链条比单独观察任意一对数据更有信息。
\end{example}

\begin{remark}
显示偏好定理的方向必须严格区分：数据满足公理，推出存在某类效用函数；数据违反公理，推出不存在满足相应性质的理性化。不能把“某个简单效用函数拟合失败”误写成“所有效用函数都无法理性化”。
\end{remark}

\subsection{GARP 不等式链的完整展开}
令第 $t$ 期观察到的选择为 $x^t$，价格为 $p^t$，收入为 $m^t=p^t\cdot x^t$。若 $x^s$ 在第 $t$ 期可负担，则
\[
 p^t\cdot x^s\le p^t\cdot x^t
 \quad\Longrightarrow\quad x^t\succeq x^s.
\]
若沿若干期得到 $x^1\succeq x^2\succeq\cdots\succeq x^k$，且最后一期严格显示排除 $x^1$，则理性化要求
\[
 u(x^1)\ge u(x^2)\ge\cdots\ge u(x^k)>u(x^1),
\]
矛盾。这个推导将图形中的“可负担边”转化为效用数值上的不等式链；GARP 的核心不是某一条价格比较，而是所有可负担边的传递闭包。

\begin{example}[三期链的逐项计算]
假设 $p^1\cdot x^2\le p^1\cdot x^1$、$p^2\cdot x^3\le p^2\cdot x^2$，且
$p^3\cdot x^1<p^3\cdot x^3$。前两式分别推出 $u(x^1)\ge u(x^2)$ 与
$u(x^2)\ge u(x^3)$；若数据还含有从 $x^3$ 回到 $x^1$ 的显示弱偏好链，便得到严格闭环。故检查时必须把“直接边”和“闭包边”分开记录，避免把一条不可负担关系误读为偏好关系。
\end{example}

\begin{proposition}[尺度变换不改变显示偏好]
若每期价格与收入同时乘以同一 $a_t>0$，则预算可负担关系及 WARP、SARP、GARP 的判定均不变。
\end{proposition}
\begin{proof}
$(a_tp^t)\cdot x^s\le a_tm^t$ 当且仅当 $p^t\cdot x^s\le m^t$。因此每一条可负担边保持不变，由这些边生成的传递闭包也保持不变，所有环和严格边的存在性自然不变。
\end{proof}

\subsection{有限数据的手算程序与逻辑检查}
对有限选择数据，理性化检验可以按四列进行：第一列记录观察期 $t$；第二列记录被选择束 $x^t$；第三列计算每个其他束在该期的支出 $p^t\cdot x^s$；第四列把可负担关系转化为有向边 $t\to s$。完成表格后，先查两期反向边，再查传递闭包中的环，最后检查环上是否至少存在一条严格边。没有严格边的闭环只说明无差异，不构成 GARP 违反；有严格边的闭环才推出效用不等式矛盾。

设三期数据的支出矩阵为 $E_{ts}=p^t\cdot x^s$。对角线 $E_{tt}=m^t$ 是实际支出。若 $E_{ts}\le E_{tt}$，就在矩阵中标记 $t\to s$；若不等号严格且 $x^s\ne x^t$，再标记为严格边。于是，GARP 的计算对象不是价格矩阵本身，而是由 $E$ 的逐行比较生成的标记矩阵。

\begin{example}[四期闭包的检查顺序]
若直接边为 $1\to2$、$2\to3$、$3\to4$，算法先生成传递边 $1\to3$、$1\to4$、$2\to4$。若另有 $4\to1$，则出现闭环；此时只需检查四条边中是否有一条严格边。若有，便得到 $u(x^1)\ge u(x^2)\ge u(x^3)\ge u(x^4)>u(x^1)$；若没有，四个束可以属于同一无差异类。将“闭包”与“严格性”分开，是避免把普通循环误判为违反的关键。
\end{example}

\subsection{从 WARP 到 SARP 再到 GARP}
WARP 只检查两期直接反转，SARP 检查直接边与传递边形成的严格循环，GARP 则允许无差异并只排除带严格边的闭环。因此三者的差别不在于是否使用价格，而在于所允许的闭包范围和对并列关系的处理方式。两期数据中不存在更长路径，故 SARP 的检验退化为 WARP；三期及以上时，单独检查每一对边不再足够。

若数据满足 GARP，Afriat 不等式给出 $U_t\le U_s+\lambda_s p^s\cdot(x^t-x^s)$，其中 $\lambda_s>0$。逐项解出 $U_t$ 与 $\lambda_s$ 后，可构造分段线性凹效用 $u(x)=\min_s\{U_s+\lambda_s p^s\cdot(x-x^s)\}$。在观测点 $x^t$ 处，第 $t$ 个平面取值为 $U_t$，其余平面不低于 $U_t$，所以 $u(x^t)=U_t$；在预算内任意束都不会超过该期观测效用。

\subsection{三期数据的逐项理性化核对}
设三期观测为 $(p^t,x^t)$，并假定每期选择都耗尽预算。对任意两期 $s,t$，先计算
\[
 p^s\cdot x^t-p^s\cdot x^s.
\]
若该差值不大于零，则 $x^t$ 在第 $s$ 期可负担，观察到的选择意味着
$x^s\succeq x^t$；若差值严格小于零且 $x^t\ne x^s$，则显示偏好为严格关系。
因此，理性化检验不能只看价格向量是否变化，而必须逐期计算“另一期选择在本期预算中的支出”。

例如，若
\[
 p^1=(1,2),\quad x^1=(2,1),\qquad
 p^2=(2,1),\quad x^2=(1,2),
\]
则
\[
 p^1\cdot x^1=4,\quad p^1\cdot x^2=5;
 \qquad
 p^2\cdot x^2=4,\quad p^2\cdot x^1=5.
\]
两期互不显示可负担，故不存在由这两期单独构成的直接环。若第三期选择
$x^3$ 同时落入前两期预算约束，就必须把三条弱不等式连起来检查；一旦得到
$u(x^1)\ge u(x^2)\ge u(x^3)>u(x^1)$，矛盾来自闭环，而不是来自任意一条单独比较。

\begin{proposition}[尺度变换不改变显示偏好]
若每期价格与收入同时乘以同一 $a_t>0$，则预算可负担关系及 WARP、SARP、GARP 的判定均不变。
\end{proposition}
\begin{proof}
$(a_sp^s)\cdot x^t\le a_sm^s$ 当且仅当 $p^s\cdot x^t\le m^s$。因此每一条可负担边都保持不变，
由这些边生成的传递闭包也保持不变，所有环和严格边的存在性自然不变。
\end{proof}
